% !TEX program = xelatex
\documentclass[UTF8,a4paper,12pt]{ctexart}
\usepackage[margin=2.5cm]{geometry}
\usepackage{amsmath,amssymb,booktabs,graphicx,siunitx,tikz,listings,xcolor}
\usepackage[hidelinks]{hyperref}
\usepackage{fancyhdr}
\usepackage{enumitem}

\pagestyle{fancy}
\fancyhf{}
\fancyhead[L]{学科竞赛论文示例}
\fancyhead[R]{\thepage}
\setlength{\headheight}{15pt}
\setlength{\parindent}{2em}
\setlength{\parskip}{0.3em}
\sisetup{per-mode=symbol}
\lstset{basicstyle=\ttfamily\small,frame=single,breaklines=true,columns=fullflexible}

\newcommand{\ProjectTitle}{低压能量回馈变流器控制系统设计}
\newcommand{\TeamName}{示例队伍}
\newcommand{\University}{示例大学}
\newcommand{\Competition}{电子设计竞赛}

\begin{document}

\begin{titlepage}
  \centering
  {\zihao{2}\bfseries \Competition 参赛论文\par}
  \vspace{2.2cm}
  {\zihao{1}\bfseries \ProjectTitle\par}
  \vspace{1.5cm}
  {\zihao{4}队伍：\TeamName\par}
  {\zihao{4}单位：\University\par}
  \vfill
  {\zihao{4}\today\par}
\end{titlepage}

\begin{abstract}
本文给出一套面向低压、隔离实验平台的能量回馈变流器控制系统设计示例。论文结构覆盖需求拆解、总体方案、硬件与软件设计、控制算法、保护策略、测试方法、结果分析和成本评估。示例数据仅用于演示排版，参赛时必须替换为真实测量结果。

\textbf{关键词：} 能量回馈；变流器；嵌入式控制；故障保护；实验验证
\end{abstract}

\tableofcontents
\newpage

\section{任务要求与指标分解}
将题目要求转换为可测量指标，并为每项指标给出工况、仪器、计算方法和通过标准。表\ref{tab:requirements}给出示例。

\begin{table}[htbp]
\centering
\caption{指标分解示例}
\label{tab:requirements}
\begin{tabular}{llll}
\toprule
指标 & 目标值 & 测量方法 & 状态 \\
\midrule
输出电流误差 & $\leq\SI{2}{\percent}$ & 四线法采样与对比 & 待实测 \\
稳态效率 & $\geq\SI{85}{\percent}$ & 输入/输出功率计算 & 待实测 \\
保护关断时间 & $\leq\SI{10}{\milli\second}$ & 示波器单次触发 & 待实测 \\
\bottomrule
\end{tabular}
\end{table}

\section{总体方案}
系统由功率变换、采样调理、数字控制、通信上位机和独立保护链路组成。图\ref{fig:architecture}使用 TikZ 绘制，无需外部图片文件。

\begin{figure}[htbp]
\centering
\begin{tikzpicture}[node distance=1.6cm,>=latex,every node/.style={draw,rounded corners,align=center,minimum width=2.7cm,minimum height=0.9cm}]
\node (power) {低压功率级};
\node (sample) [right of=power,xshift=2.2cm] {电压/电流采样};
\node (mcu) [right of=sample,xshift=2.2cm] {STM32 控制器};
\node (pc) [below of=mcu] {RS485 上位机};
\node (protect) [below of=sample] {硬件保护链路};
\draw[->] (power) -- (sample);
\draw[->] (sample) -- (mcu);
\draw[<->] (mcu) -- (pc);
\draw[->] (protect) -- (power);
\draw[->] (sample) -- (protect);
\end{tikzpicture}
\caption{系统总体架构}
\label{fig:architecture}
\end{figure}

\section{控制算法}
电流环采用离散 PI 控制。误差与控制量分别为
\begin{align}
  e[k] &= i^*[k]-i[k],\\
  u[k] &= K_p e[k] + K_i T_s\sum_{j=0}^{k} e[j].
\end{align}
为防止积分饱和，应在输出限幅后执行条件积分或反算抗饱和，并对参考值增加斜率限制。

\section{安全与故障处理}
\begin{enumerate}[label=\arabic*)]
  \item 上电、复位、通信超时和故障状态均保持 PWM 为零；
  \item 急停与硬件过流关断不依赖软件；
  \item 故障锁存后仅允许在停机且测量恢复正常时清除；
  \item 先完成无功率、低压限流测试，再逐级提高工况。
\end{enumerate}

\section{测试方法与结果}
所有结果应记录仪器型号、环境、固件版本、原始 CSV 和不确定度。不得只保留截图。示例结果见表\ref{tab:results}。

\begin{table}[htbp]
\centering
\caption{测试结果示例（非真实竞赛数据）}
\label{tab:results}
\begin{tabular}{lrrr}
\toprule
工况 & 参考电流/A & 实测电流/A & 相对误差/\% \\
\midrule
低负载 & 0.50 & 0.49 & 2.00 \\
中负载 & 1.00 & 0.99 & 1.00 \\
高负载 & 1.50 & 1.48 & 1.33 \\
\bottomrule
\end{tabular}
\end{table}

\section{结论}
本模板给出了竞赛论文的完整技术叙事骨架。正式提交前应删除所有示例数据，并保证论文、代码、原理图、PPT 与实物参数一致。

\begin{thebibliography}{9}
\bibitem{ref1} 示例作者. 功率电子系统设计方法[M]. 示例出版社, 2026.
\bibitem{ref2} STMicroelectronics. STM32F103 Reference Manual[Z].
\end{thebibliography}

\end{document}
